% =============================================================================
\section{Résultats et comparaisons}
% =============================================================================

\begin{frame}{Statique vs hybride : la hiérarchie confirmée}
  \footnotesize
  \begin{columns}[T,onlytextwidth]
    \begin{column}{0.52\textwidth}
      \begin{block}{À budget temps égal}
        \begin{itemize}
          \item Recuit simulé $\approx$ Tabou $\times$ Recuit (tête) ;
          \item puis Tabou $\approx$ Génétique $\times$ Tabou ;
          \item SMA au milieu, freinés par le coût de coordination ;
          \item ACO et génétique distancés, surtout à l'échelle réelle.
        \end{itemize}
      \end{block}
      \begin{exampleblock}{Le message}
        \footnotesize Les \textbf{hybrides} apportent un gain net
        (Génétique $\times$ Tabou, Fourmis $\times$ Tabou) ; les méthodes
        \textbf{locales} restent les plus fiables sur un paysage rugueux.
      \end{exampleblock}
    \end{column}
    \begin{column}{0.46\textwidth}
      \centering
      \includegraphics[width=\linewidth,height=0.28\textheight,keepaspectratio]{grande_boxplot.png}\\[2pt]
      {\tiny\color{mcmPrimary}\textbf{Dispersion} de la fitness finale
        (grande échelle).}\\[4pt]
      \begin{alertblock}{Fiabilité}
        \footnotesize 10/10 méthodes atteignent l'optimum exact en petite
        instance ; aucune ne viole les contraintes dures.
      \end{alertblock}
    \end{column}
  \end{columns}
  \source{Synthèse des campagnes 20 graines : petite (7), grande (300) et réelle
    (582) instances.}
\end{frame}

\begin{frame}{Résultats : prédictions, oracle et capacité}
  \small
  \centering
  \renewcommand{\arraystretch}{1.40}
  \setlength{\tabcolsep}{11pt}
  \begin{tabular}{lrrr}
    \toprule
    Capacité & Médianes & Modèle & Oracle \\
    \midrule
    42 lits & 246 & \textbf{251} & 260 \\
    10 lits & \textbf{111} & 108 & 115 \\
    \bottomrule
  \end{tabular}
  \vspace{0.38cm}
  \begin{columns}[T,onlytextwidth]
    \begin{column}{0.49\textwidth}
      \begin{block}{42 lits : non saturés}
        \small Au plus 30 lits occupés en fin de journée.
        Le modèle gagne 5 cas sur les médianes, mais reste
        9 cas derrière l'oracle.
      \end{block}
    \end{column}
    \begin{column}{0.49\textwidth}
      \begin{alertblock}{10 lits : saturation}
        \small Les 10 lits sont occupés en fin de journée
        sur \textbf{21 des 22 jours} opératoires.
      \end{alertblock}
    \end{column}
  \end{columns}
  \source{Cas terminés sur 279 ; politique réactive, fermeture 10h--12h, 2 salles, graine 0. Oracle = durées futures révélées, pas optimum garanti.}
\end{frame}

\begin{frame}{Résultats : mieux chercher, mieux informer}
  \footnotesize
\begin{columns}[T,onlytextwidth]
\begin{column}{.48\textwidth}
\textbf{Choix de la méthode : budget identique}\par
300 patients / 80 vacations ; 3 s ; 20 graines.
\par\vspace{0.8mm}
\centering
\renewcommand{\arraystretch}{1.1}
\begin{tabular}{@{}lr@{}}
\toprule
Méthode & Écart moyen\\\midrule
\textbf{Recuit simulé} & \textbf{0,015}\\
\textbf{Tabou $\times$ Recuit} & \textbf{0,018}\\
Tabou & 0,240\\
Génétique $\times$ Tabou & 0,261\\\bottomrule
\end{tabular}
\par\vspace{0.5mm}
\raggedright
Écart à la meilleure solution observée,\newline pas à un optimum prouvé.
\par\vspace{0.5mm}
\textbf{582 / 280 :} même groupe de tête ; meilleure solution faisable, lits à 42/42.
\end{column}
\begin{column}{.48\textwidth}
\textbf{Qualité de l'information : mêmes ressources}\par
279 interventions ; 42 lits ; fermeture ; réactif.
\par\vspace{0.8mm}
\centering
\renewcommand{\arraystretch}{1.1}
\begin{tabular}{@{}lr@{}}
\toprule
Information & Terminées / 279\\\midrule
Médianes & 246\\
\textbf{Prédictions ML} & \textbf{251}\\
Oracle & 260\\\bottomrule
\end{tabular}
\par\vspace{0.5mm}
\raggedright
\textbf{ML : +5 interventions} par rapport aux médianes ; 9 de moins que l'oracle.
\par\vspace{0.5mm}
À \textbf{10 lits}, médianes 111 / ML 108 : le gain dépend des ressources.
\end{column}
\end{columns}
\vspace{0.5mm}
\centering
{\color{McmTerracotta}\bfseries Recuit et Tabou $\times$ Recuit en tête ;
la prédiction aide, sans gain systématique.}
\end{frame}

\begin{frame}{Aléas : deux compromis opérationnels}
  \footnotesize
\begin{columns}[T,onlytextwidth]
\begin{column}{.45\textwidth}
\textbf{Fermeture du bloc : adapter l'exécution}\par
279 interventions ; ML ; 42 lits ; salle fermée 10h--12h.
\par\vspace{0.8mm}
\centering
\renewcommand{\arraystretch}{1.1}
\begin{tabular}{@{}lrr@{}}
\toprule
Plan & Terminées & Dépass. (min)\\\midrule
Figé & 255 & 2\,911\\
Réactif & 251 & 2\,716\\\bottomrule
\end{tabular}
\par\vspace{0.8mm}
\raggedright
{\color{McmTerracotta}\bfseries $-195$ min de dépassement,\\mais 4 interventions de moins.}
\par\vspace{0.5mm}
Dépassement : occupation et remise en état après 17h.
\end{column}
\begin{column}{.52\textwidth}
\textbf{LOS : anticiper les séjours longs}\par
420 patients / 30 jours ; 42 lits ; Monte Carlo.
\par\vspace{0.8mm}
\centering
\renewcommand{\arraystretch}{1.1}
\setlength{\tabcolsep}{4pt}
\begin{tabular}{@{}lrr@{}}
\toprule
 & Dégradation & Coût moyen\\
 & \multicolumn{2}{c}{Ponctuel $\rightarrow$ borne haute}\\\midrule
Recuit & $109\rightarrow13$ & $5\,316\rightarrow5\,384$\\
Tabou & $115\rightarrow17$ & $5\,313\rightarrow5\,376$\\
Tabou $\times$ Recuit & $112\rightarrow6$ & $5\,311\rightarrow5\,413$\\\bottomrule
\end{tabular}
\par\vspace{0.8mm}
\raggedright
{\color{McmTerracotta}\bfseries Moins de dégradation,\\mais un coût moyen plus élevé.}
\par\vspace{0.5mm}
Dégradation : $\mathbb E[C]-C_{\mathrm{ponctuel}}$\newline pour le \textbf{même} planning.
\end{column}
\end{columns}
\vspace{0.5mm}
\centering
{\color{McmTerracotta}\bfseries La robustesse a un prix :
choisir le compromis, sans gain universel.}
\end{frame}